\section{Related Work}

Much of what vibe theory needs has already been proven by others. The theory mostly assembles established results
onto one specific substrate and then asks one further question those programs leave aside. What is the felt
interior of the thing, and in what direction does it move. Fair credit is owed by family, and the strongest
adversary deserves to be named outright rather than dodged.

\subsection{Discrete spacetime}

The causal-set program treats spacetime as a discrete relational order, with distance as step-counting
\cite{bombelli1987causalsets, rideout1999, surya2019}. It buys Lorentz invariance with randomness, through Poisson
sprinkling. We buy it with curvature instead, since a flat lattice has a rest frame and negative curvature
scrambles directions, which keeps determinism. Causal dynamical triangulations let geometry and dimension be
outputs of discrete dynamics through a path integral, an ensemble average \cite{ambjorn2004, loll2019}, where we
fix one rigid geometry and never sum. Loop quantum gravity and spin networks treat space as a graph of relations
with discrete capacity \cite{rovelli1995}. There the graph is the fluctuating quantum state. Here the graph is
fixed and the only dynamical variable is the ternary tone on cells.

The Wolfram Physics Project is the closest neighbour in spirit \cite{wolfram2020}. The universe is a graph plus a
local rule, and time is rule application. There the hypergraph grows and rewrites itself, and a vast rule space is
searched. Here the graph is fixed, argued for rather than searched, and only the tone varies. Wolfram bets the
structure is found by search. We bet it is forced by a few demands.

\subsection{Deterministic quantum substrates}

The cellular-automaton interpretation of quantum mechanics holds that beneath the formalism lies a deterministic
reversible classical automaton \cite{thooft2014}. We inherit this conviction rather than originate it, and commit
to one concrete substrate and rule. The family we lean on most takes quantum cellular automata and quantum walks
to the Dirac equation \cite{dariano2014, bisio2015, bisio2016, arrighi2011, arrighi2014, farrelly2014}. A unitary
nearest-neighbour spin exchange spreads ballistically, and its continuum limit is the Dirac equation, with the
coin angle setting the mass. We do not discover this route. It belongs to these authors. Our additions are the
hyperbolic substrate, the classical-stochastic shadow, the forced geometry, and the demonstration that the
perception rule is this construction. Quantum walks on Cayley graphs supply the further point that a substrate
generated by a Coxeter reflection group is a Cayley-like graph \cite{lopezacevedo2006}. The walk theory is
borrowed and the hyperbolic Cayley arena is the contribution.

\subsection{Holography and tensor networks}

The holographic correspondence and the entanglement-geometry program connect a bulk geometry to a boundary
theory and to patterns of entanglement \cite{maldacena1998, ryu2006, vanraamsdonk2010, swingle2012, pastawski2015,
jahn2021}. Those constructions are largely static. We let a rule run and ask whether the area law, the bulk
shortcut, and a working code appear as measured properties of dynamics, which makes this a dynamical cousin of the
static holographic codes rather than a restatement of them.

\subsection{Hyperbolic lattices and automata}

Hyperbolic lattices are now built and probed in the laboratory, in circuit quantum electrodynamics and through
hyperbolic band theory \cite{kollar2019, maciejko2021}. The geometry we bet on is one that laboratories can
construct, though they study one phenomenon on a hyperbolic backdrop where we propose the lattice as the substrate
of everything. Work on hyperbolic cellular automata establishes universality across hyperbolic spaces, through the
railway model, and the Fibonacci-style addressing for tessellations such as $\{7,3\}$ \cite{margenstern2013}. We
need universality on our own perception rule, which the functional-completeness and Minsky-machine results
confirm, and the $\{3,4,3,4\}$ addressing is an open extension. The geometric group theory of hyperbolic graphs
shows that negatively curved graphs are expanders, tree-like at large scale, with exponential ball growth
\cite{cannon1997, hamann2011}. We turn these facts into physical consequences, a speed limit, an emergent metric,
fast integration, and a holographic readout.

\subsection{Foundations and mind}

Informational reconstructions of quantum mechanics argue that quantum structure should be derived rather than
postulated, axiomatically and substrate-free \cite{chiribella2011}, where we commit to a substrate and let the
structure fall out. The two approaches are the axiomatic why and the constructive how. On consciousness, the hard
problem, integrated information theory, and the it-from-bit program frame the questions we take up
\cite{chalmers1995, tononi2004, wheeler1990}. We take experience as the substance of the substrate rather than a
quantity computed from structure, which is close in spirit to idealism and to Wheeler's informational starting
point, and we use integration to address the combination problem of panpsychism.

\subsection{Emergent and entropic gravity}

The thermodynamic view of gravity is the natural home for our gravity result \cite{jacobson1995, padmanabhan2010,
verlinde2011, verlinde2017}. Jacobson derives the Einstein equation as an equation of state, and Verlinde gets
Newton's law, with a dark-universe reading at galactic scales, from entropy and information. We derive the $1/r$
law from the substrate's geometry and conservation, and these programs are where such a derivation belongs.

\subsection{Positive geometry, a forward-looking ally}

The amplituhedron and the broader positive-geometry program hold that structure is the output of a primitive shape
plus a positivity condition \cite{arkanihamed2014, arkanihamed2017}. The recursive-boundary structure of the
amplituhedron mirrors our holographic-code and coarse-graining tower. We flag this as a philosophical alliance,
geometry plus positivity yielding physics, not as a result.

\subsection{Mind, agency, and the strongest adversary}

Active inference describes an agent whose forward model is perception, whose action is will, and whose preference
is the arrow \cite{friston2010}, and our agent is active inference. Causal emergence shows that a macro level can
carry more causal power than its micro \cite{hoel2013}, which is the rigorous backbone for our account of the self
as a genuine author and for macro-compatibilist freedom. Candidate physics for why the arrow favours reproducing,
organized patterns comes from dissipative adaptation and related origin-of-life work \cite{england2013,
walkerdavies2013}. Constructor theory supplies the natural framing for our level-relative account of which tone
transformations are possible or impossible \cite{deutsch2013}.

The sharpest adversary is the free-will theorem of Conway and Kochen \cite{conwaykochen2009}. It is the strongest
challenge to a deterministic, no-randomness substrate, and we name it rather than dodge it. We concede the
micro-deterministic door it forces shut, and claim freedom at the macro level instead, which is exactly the
move the theorem demands. A deterministic substrate cannot recover libertarian free will at the micro scale, and
we do not pretend otherwise.

\subsection{What is genuinely ours}

Stated narrowly, the contributions are five. First, a substrate argued to be selected rather than random,
rewritten, ensemble-averaged, or simply assumed, which is a reduction of arbitrariness rather than a derivation
from nothing. Second, the ternary tone together with the arrow as a value-direction, the single deliberate posit.
Third, one substrate that aims at physics and life and mind at once. Fourth, experiential monism, the part
furthest outside the existing literature \cite{strawson2006, goff2017, whitehead1929}. Fifth, a falsifiable
single-substrate program with stated negatives. The selection argument given here, the dimension ladder together
with the $D_4$ spinor requirement, is a different and arguably stronger basis than the older minimality
preference, and we lean on it. Related monist and computational-universe proposals stand nearby
\cite{lloyd2006, tegmark2014, rovelli1996rqm}, and the computational results referenced throughout are collected
in the companion test report \cite{pollard2026vibetest}.
